% Source: https://artofproblemsolving.com/community/c3h1656531_2018_ms_mte_sprint_all
% Stable identity: aops-topic-1656531-q18
The units digit of $m^3$ depends only on the units digit of $m$: for $m\equiv1,2,\dots,9,0\pmod{10}$ it is $1,8,7,4,5,6,3,2,9,0$, which sum to $45$. Hence the running sums of these units digits for $n=1,\dots,20$ are
\[
1,9,16,20,25,31,34,36,45,45,\ 46,54,61,65,70,76,79,81,90,90,
\]
and after that they increase by $90$ every $20$ steps, so their units digits $f(n)$ repeat with period $20$:
\[
f(1),\dots,f(20)=1,9,6,0,5,1,4,6,5,5,\ 6,4,1,5,0,6,9,1,0,0 .
\]
One period sums to $74$, and the partial sums within a period are
\[
1,10,16,16,21,22,26,32,37,42,48,52,53,58,58,64,73,74,74,74 .
\]
So the possible values of $f(1)+\cdots+f(k)$ are exactly the numbers $74q+s$ with $q\ge0$ and $s$ in this list, i.e.\ the residues modulo $74$ that occur are
\[
0,1,10,16,21,22,26,32,37,42,48,52,53,58,64,73 .
\]
Now $70\equiv70$, $158\equiv10$, $449\equiv5$, $1615\equiv61$ and $7904\equiv60\pmod{74}$. Only $158=2\cdot74+10$ is possible (it is the value for $k=42$). The answer is $\boxed{\textbf{(b)}\ 158}$.
